
% ---------------------------------------------------------------------------
% Canonical addition for Inquiry Calculus v1.1
% Non-constitutional explanatory/derived section.
% ---------------------------------------------------------------------------

\section{Canonical Example: Paired Actuality Spine, Residual Cues, and Recursive Inquiry}
\label{sec:paired-actuality-spine}

\begin{remark}[Status of this section]
This section introduces no new constitutional primitive, typing rule, authority
class, history species, or semantic law.  It is a canonical worked
interpretation of structures already present in the calculus: typed relations,
partial binding, holes and residuals, reification, probe occurrence,
arrangement/succession separation, operator composition, folding, reopening,
and protected regeneration.

Its purpose is to make explicit an especially important fact about inquiry:
the realized inquiry history itself is an ordinary relational object.  It may
therefore be questioned, partially bound, have components removed and
regenerated, be compared across occurrences, abstracted, folded into methods,
reopened, and used to construct new objects of inquiry.
\end{remark}

\subsection{One realized inquiry occurrence}

Let \(A_t\) denote the represented arrangement available at an inquiry step,
let \(q_t\) denote the realized question/probe/program presented at that step,
let \(r_t\) denote the realized return, and let \(A_{t+1}\) denote the
successor arrangement.

A realized occurrence may be viewed as the already-defined probe event
\[
\boxed{
E_t=(A_t,q_t,r_t,A_{t+1})
}
\]
or equivalently
\[
\boxed{
A_t\xrightarrow[q_t]{r_t}A_{t+1}.
}
\]

This notation does not identify arrangement with succession.  It displays
their realized coupling while retaining their independent projections:
\[
\pi^-_A(E_t)=A_t,
\qquad
\pi_Q(E_t)=q_t,
\qquad
\pi_R(E_t)=r_t,
\qquad
\pi^+_A(E_t)=A_{t+1}.
\]

Thus a realized history
\[
H=E_0E_1E_2\cdots
\]
admits, among many others, the projections
\[
H_A=\pi_A(H),
\qquad
H_Q=\pi_Q(H),
\qquad
H_R=\pi_R(H),
\qquad
H_{QR}=\pi_{Q,R}(H).
\]

These are views of one authoritative occurrence history, not independent
authoritative stores.

\subsection{Two reciprocal trace projections and their coupled whole}

The inquiry trace may be projected as an alternating relation
\[
\boxed{
Q\xrightarrow{\alpha}R\xrightarrow{\kappa}Q.
}
\]

Here \(\alpha\) is the realized-return relation induced by asking/probing under
the applicable arrangement and binding, while \(\kappa\) is the
return-dependent continuation relation by which a realized return makes a
subsequent question available or consequential.

Suppressing arrangement only for the purpose of displaying the relational
shape,
\[
\alpha(q_t,r_t)
\]
and
\[
\kappa(r_t,q_{t+1})
\]
give the alternating trace
\[
q_0\xrightarrow{\alpha}r_0
\xrightarrow{\kappa}q_1
\xrightarrow{\alpha}r_1
\xrightarrow{\kappa}q_2\cdots.
\]

The question-to-question projection is the relational composite
\[
\boxed{
T_Q=\kappa\circ\alpha:Q\rightsquigarrow Q,
}
\]
so that
\[
T_Q(q_t,q_{t+1})
\iff
\exists r.\,
\alpha(q_t,r)\land\kappa(r,q_{t+1}).
\]

Dually, the return-to-return projection is
\[
\boxed{
T_R=\alpha\circ\kappa:R\rightsquigarrow R,
}
\]
so that
\[
T_R(r_t,r_{t+1})
\iff
\exists q.\,
\kappa(r_t,q)\land\alpha(q,r_{t+1}).
\]

Hence there are two useful projections and their coupled whole:
\[
\boxed{
\text{question trace},
\qquad
\text{return trace},
\qquad
\text{alternating question--return trace}.
}
\]

The first relates questions through returns; the second relates returns
through questions; the alternating trace retains the mediator that each
projection hides.  This is ordinary relational composition and hiding, not a
new ontology.

\subsection{Human and LLM bindings}

The same relational shape may be instantiated differently.

For an LLM binding,
\[
q_t
\]
may be a compiled natural-language question or inquiry program,
\[
r_t
\]
the raw model return, and
\[
A_t
\]
the active/contextual arrangement at that occurrence.  Fixed model parameters
\(\Theta\) parameterize a possibility relation
\[
G_\Theta(A_t,q_t)\rightsquigarrow R,
\]
while the authoritative occurrence history records the particular realized
return:
\[
r_t\in G_\Theta(A_t,q_t).
\]

The realized pair \((q_t,r_t)\) belongs to history.  Ordinary inference does
not thereby write that occurrence back into the fixed weights \(\Theta\).

For a human binding, \(q_t\) need not be verbal.  It may be a perceptual
orientation, bodily movement, manipulation, experiment, deliberate question,
or other probe; \(r_t\) is the corresponding perceptual/environmental return
at the chosen grain; and \(A_t\) is the relevant bodily, environmental, and
represented arrangement.  Physical/spatiotemporal succession is supplied by
the binding rather than by the abstract calculus.

Thus both bindings instantiate the same derived schema
\[
\boxed{
A_t\xrightarrow[q_t]{r_t}A_{t+1}.
}
\]

\subsection{Actuality and possibility}

The paired spine records actuality; it does not exhaust possibility.

For an LLM binding,
\[
r_t\in G_\Theta(A_t,q_t)
\]
states that the realized return belonged to the model's admissible conditional
possibility field at that occurrence.  In general,
\[
G_\Theta(A_t,q_t)\neq\{r_t\}.
\]

More generally, the calculus distinguishes a reasoner's represented
possibility from the binding's admitted actuality.  A realized return may
force reopening when it was not represented in the prior live possibility
field.  Therefore the safe relation is:
\[
\boxed{
\text{realized actuality is one realized path through an admitted possibility
relation, while the reasoner's prior represented possibility may be
incomplete.}
}
\]

The possibility field need not be stored as an explicit enumeration.  It may
be generated compositionally from previously acquired relational structure,
operators, and representations.

\subsection{The inquiry history is itself an object of inquiry}

Because questions, returns, events, paths, and arrangements are reifiable
forms, any projection or composition of the history may itself enter a new
relation.

For example:
\[
\ulcorner q_i\urcorner,
\qquad
\ulcorner r_i\urcorner,
\qquad
\ulcorner(q_i,r_i)\urcorner,
\qquad
\ulcorner T_Q\urcorner,
\qquad
\ulcorner T_R\urcorner
\]
may be compared, distinguished, partially bound, or placed on either side of
a further distinction.

Thus one may lawfully ask:
\[
?\rho[\rho(q_i,q_j)],
\]
\[
?\rho[\rho(r_i,r_j)],
\]
or
\[
?\Delta[
\Delta(
\ulcorner(q_i,r_i)\urcorner,
\ulcorner(q_j,r_j)\urcorner)
].
\]

The sixfold reciprocal construction, section/fiber construction, breaker
search, representation invention, and protected-equivalence machinery apply
unchanged.  Inquiry therefore has no privileged exemption from inquiry.

\subsection{Removing a return: the residual cue}

A particularly useful derived construction arises by removing a return while
retaining the relations on either side of it.

Given the local alternating trace
\[
q_t\xrightarrow{\alpha}r_t\xrightarrow{\kappa}q_{t+1},
\]
remove \(r_t\) but preserve the two relations:
\[
q_t\xrightarrow{\alpha}\square\xrightarrow{\kappa}q_{t+1}.
\]

The lawful fillers are
\[
\boxed{
\operatorname{Cue}_R(q_t,q_{t+1})
=
\{
r:R:
\alpha(q_t,r)\land\kappa(r,q_{t+1})
\}.
}
\]

This is not a new primitive ``cue'' type.  It is exactly the solution fiber
of the residual hole:
\[
\operatorname{Cue}_R(q_t,q_{t+1})
=
\operatorname{Sol}_{\alpha,\kappa}(q_t,q_{t+1}).
\]

If
\[
\operatorname{Cue}_R(q_t,q_{t+1})/
{\equiv_{\mathcal H}}
=
\{[r_t]_{\mathcal H}\},
\]
then the omitted return is regenerable, to the protected grain, from the
surrounding inquiry relation.

If several protected-distinct returns remain,
\[
\left|
\operatorname{Cue}_R(q_t,q_{t+1})/
{\equiv_{\mathcal H}}
\right|>1,
\]
the residual is not a complete reconstruction.  Its ambiguity is itself a
lawful object of inquiry and may generate a separator question.

\subsection{Removing a question: reciprocal cue construction}

The same operation applies in the reciprocal direction.

Given
\[
r_{t-1}\xrightarrow{\kappa}q_t\xrightarrow{\alpha}r_t,
\]
remove \(q_t\):
\[
r_{t-1}\xrightarrow{\kappa}\square\xrightarrow{\alpha}r_t.
\]

The lawful question fillers are
\[
\boxed{
\operatorname{Cue}_Q(r_{t-1},r_t)
=
\{
q:Q:
\kappa(r_{t-1},q)\land\alpha(q,r_t)
\}.
}
\]

Thus differences between retained return-side cues may regenerate questions,
just as question-side structure may regenerate omitted returns.

Again, this is only the ordinary hole/fiber construction applied to the
inquiry trace.

\subsection{Cue-based reconstruction as a memory strategy}

When direct retention or retrieval of detailed returns is expensive or
unreliable, a reasoner may preserve a smaller structure that retains the
easily reproducible or compressible inquiry path together with enough
residual constraints to regenerate omitted material.

Schematically,
\[
q_t
\longrightarrow
\operatorname{Cue}_R(q_t,q_{t+1})
\longrightarrow
q_{t+1}.
\]

Such a cue is useful exactly to the extent that it reduces the protected
solution fiber.  It need not literally reconstruct the historical wording of
the omitted response.  What matters is protected regenerative sufficiency:
\[
\boxed{
\operatorname{Regen}_{\mathcal H}
\left(
(q_t,\operatorname{Cue}_R,q_{t+1}),
r_t
\right).
}
\]

A human reasoner with weak verbatim retention may therefore reconstruct
content from relational position rather than from direct episodic recall.
This is a possible derived strategy, not a universal claim about human memory.

Likewise, a compact LLM-facing memory system may retain recurrent question
programs and references into authoritative history while reopening detailed
arrangements or returns only when the current inquiry requires them.

\subsection{Program-stream compression}

The arrangement stream may be large:
\[
A_0,A_1,A_2,\ldots.
\]

The question/program stream
\[
q_0,q_1,q_2,\ldots
\]
is often much smaller and may contain recurrent typed paths.

If
\[
P=q_n\circ\cdots\circ q_1
\]
recurs under comparable contracts and protected continuation behavior, the
existing method-fold machinery may promote
\[
\boxed{
m\doteq_{\mathcal H}P.
}
\]

The fold does not delete its supporting occurrences or arrangements.  It
stores references to them and remains subject to its applicability,
provenance, continuation-descent, recovery, and unlock contracts.

Accordingly, a useful implementation heuristic is:
\[
\boxed{
\textbf{COMPRESS THE REUSABLE INQUIRY PATH BEFORE ATTEMPTING
DESTRUCTIVE COMPRESSION OF THE ARRANGEMENT IT TRAVERSED.}
}
\]

This is an architectural heuristic derived from the existing compression
laws, not an additional semantic requirement.

\subsection{Traversing the memory spine}

A compact method or question-pattern graph can serve as an index into the
larger authoritative occurrence history.

A learned program node \(m\) retains references to occurrences
\[
E_{i_1},E_{i_2},\ldots,E_{i_k}.
\]

Following one such reference recovers a location on the authoritative spine:
\[
m\leadsto E_i\leadsto A_i.
\]

From that location, ordinary ancestry and succession relations permit
movement to neighboring occurrences and arrangements:
\[
\cdots
E_{i-1}
\to
A_i
\xrightarrow[q_i]{r_i}
A_{i+1}
\to
E_{i+1}
\cdots.
\]

Hence compact operator structure and large retained arrangement serve
different roles:
\[
\boxed{
\begin{aligned}
\text{operator/program graph} &: \text{compact routing and regeneration},\\
\text{actuality spine} &: \text{authoritative realized arrangement and return history}.
\end{aligned}
}
\]

The first indexes, folds, and navigates the second; it does not replace it.

\subsection{Pattern extraction is ordinary inquiry}

Patterns in the inquiry trace require no special pattern primitive.

Given two question-path fragments \(P_i,P_j\), form a distinction
\[
z=\Delta(P_i,P_j)
\]
and apply the ordinary reciprocal, variation, breaker, and protected
equivalence machinery.

Given repeated prompt/return occurrences, remove a component and inspect its
solution fiber.  If the same protected filler class is regenerated across
many contexts, a candidate invariant has been exposed.  If the filler field
splits under some context, that context is a separator and may refine the
representation.

Accordingly:
\[
\boxed{
\text{extract pattern}
=
\text{form a relation}
+
\text{vary or remove components}
+
\text{inspect the surviving residual}
+
\text{test protected discrimination}.
}
\]

Compression is the inverse-direction question:
\[
\boxed{
\text{what smaller form regenerates the same protected relational structure?}
}
\]

Expansion/reopening asks:
\[
\boxed{
\text{what hidden distinction has become inspectable and therefore must be
regenerated from the fold?}
}
\]

\subsection{Why this example is central but not foundational}

The paired actuality spine is not an additional foundation beneath the
calculus.  It is a particularly transparent instantiation of the calculus's
existing closure.

Every component is already available:
\[
\begin{aligned}
&\text{question/probe} &&\text{as an open relation and executable operator},\\
&\text{return} &&\text{as preserved actuality},\\
&\text{arrangement} &&\text{as represented relational structure},\\
&\text{succession} &&\text{as an independent relation},\\
&\text{trace projection} &&\text{as relational composition/hiding},\\
&\text{cue} &&\text{as a residual solution fiber},\\
&\text{pattern} &&\text{as a reified relation under inquiry},\\
&\text{method} &&\text{as a protected fold of a recurrent typed path},\\
&\text{memory traversal} &&\text{as navigation through linked authoritative ancestry},\\
&\text{reconstruction} &&\text{as protected regenerative sufficiency}.
\end{aligned}
\]

The example therefore illustrates a general closure principle:
\[
\boxed{
\begin{gathered}
\textbf{ANY REALIZED INQUIRY TRACE CAN ITSELF BE TURNED INTO AN}\\
\textbf{OBJECT OF INQUIRY; ITS COMPONENTS MAY BE RELATED, REMOVED,}\\
\textbf{REGENERATED, DISTINGUISHED, FOLDED, REOPENED, AND}\\
\textbf{RECOMPOSED BY THE SAME CALCULUS THAT GENERATED THE TRACE.}
\end{gathered}
}
\]

This is why inquiry remains generative without requiring a separate
metacognitive architecture.  Once a relation is represented, it is available
to the same typed operations as every other represented form.
